\chapter{The notebooks: the other half of the same working life}

\section{Four of ninety}

The Whitney counts ninety notebooks of Josephine Hopper's in the Sanborn Hopper Archive and has digitised four. The Provincetown Art Association holds twenty-two diaries of 1933 to 1956, and Madeleine Larson's typescript of them, in five parts, is used here through the edition's timeline. What follows is written from the four the Whitney serves as photographs, because those are the four whose reading can be checked against the hand, and it uses the typescript where the ledgers and the notebooks need a date or a sentence that only the diary gives.

The four are not of one kind. \emph{Garrulities + Grouch}, subtitled \q{When Artists Meet}, is a diary of the marriage as it bore on her painting, kept at long intervals from 1927 to 1959. \emph{Re: 3 Wash Sq} and \emph{Battle of Wash Sq.} are the working file and the memorandum book of the fight of 1947. The black two-ring notebook is a book of lists. Between them they supply what the ledgers leave out: what the studio cost to run, what the frames were, what the pictures were called before they were named, what she thought of him, and what she thought of herself.

\section{Garrulities + Grouch}

The title page is in pencil: \q{Garrulities + Grouch.}, a rubbed-out line, \q{When Artists Meet.}, and three more erased lines nobody can read (\nb{Garrulities}{1}{97409}). The book turns to ink on page 4 and stays in ink until a blue-grey entry of 1956. Its dates are in the margin, and it was not written front to back: an entry of 25 April 1928 follows one of the 27th, and a January 1952 entry follows February's. The transcription keeps the book's order. The subject, sustained over thirty-one years and twenty-two dated entries, is stated in the second: \q{If ever E.H. comes across these notes, he wo'nt find anything that he has'nt heard several times to his face. We say everything - at least I do - + bless his heart, he takes it.} (20 Feb. 1928, \nb{Garrulities}{14}{96263}).

The longest entry, of 6 February 1928, runs ten pages and contains the grievance whole. Before the marriage he had made her eleven frames for her exhibition, carried her watercolours to the Brooklyn International, written to Chicago on her behalf. A year after it, \q{E.H. introduced the man from the Cleveland Mus. to his wife -- + never mentioned the fact that she too was an artist -! That did something to me from which I've never recovered. \ldots\ For 3 1/2 yrs. my water colors have roosted on their rack and no living mortal passing our threshold has had the vaguest curiosity regarding the nature of them.} (\nb{Garrulities}{11}{97273}). \q{for so long I have'nt had any identity at all. People have been lovely + cordial to me but only as his wife} (p. 5). \q{his contributions to his wife have been chiefly inhibitions.} (p. 12). The same entry describes his work in a sentence that is also a description of her own: \q{remote + chilled + realistic - and strong + honest - + pessimistic + succinct - as no woman's work can ever be.} (p. 5, ref 96067). And it describes him: \q{Tall + lanky + “a Lincoln-like”. Same height as Lincoln 6 ft. 4. And honest + humorous + essentially kindly. (Except when his Egotism shuts it out)} (p. 9), with the Rehn show of 1924 remembered as \q{“the best thing since Winslow Homer” was the shout at his first show at Rehns - off in a little back rooms, 11 water colors}.

The entry also contains a staged exhibition. The morning after a quarrel she laid out on his worktable four \q{Exhibits}, her fragile treasures each paired with one of his tools, an earring with a hammer, a lace handkerchief with a lump of coal, under a sign: \q{“Opportunity Gallery / Humor + destructiveness of Edward invited.” I lost a minute - could'nt think - it was Egotism - not destructiveness that I'd planned.} He put everything back tenderly and corrected her spelling of \q{Oppertunity} (p. 6, ref 97505). It is the one exhibition in the archive that she curated, and it was reviewed by its only visitor.

Then the entries that make the notebook painful to read. 25 April 1928: \q{An encounter with an Egotist is like beating one's head against a stone wall. The wall is in no wise affected.} (p. 18); \q{he beat me well over the head with his fists, or palms - I do'nt know which --- for good + sufficient reason: I'd upset his bureau drawers \ldots\ But what can you do when they're 6 ft. 4 + have hands like iron clamps} (p. 21, ref 97465); \q{I ca'nt paint out of disgust. It's a knife in the breast. He knows this --- + he works it.} (p. 23). 23 November 1930: the easel she brought into the studio in 1925 over his objection, and \q{now that his resistance is withdrawn my so slight a giftie has given up me --- died in the squabble. \ldots\ All the dishes that have been washed, the bacon fried, floors swept, socks darned --- avail nothing.} (p. 28, ref 97196). 17 February 1934, written in bed with a cold: \q{By any chance have I married a sadist ?} (p. 37); \q{I am a woman --- + deny, stamp it out having first uprooted --- I am also an artist.} (p. 43, ref 96250); and, at a Whitney party, a title offered for a picture: \q{“What'll I call it ?” I was asked + said : “Damned Egotists.” E.H. here at my side said Why not “God Damned Egotists” --- + brought down the house.} (p. 36).

The record is not one-sided, and its keeper knew it. The broken ankle of January 1931 makes him \q{King of beast and best nurse + cook + house maid + janitor of all the best artists in America to-day} (p. 31, ref 96245), and \q{I might better be writing a book about the cute things he says. \ldots\ He says it's always too intimate, no one would get it but me.} (p. 30). The wedding day, told in September 1939 because Guy Pène du Bois had told it wrong, has the Little Church Around the Corner refusing them because he had never been baptised, church after church in the July heat, and the little French church on West Fifteenth Street with the concierge and his wife as witnesses; \q{when E.H., who tells his business to no one, short: “We're going to be married!” said he.} (p. 48, ref 96072). The head margin of that page carries, fifteen years after the event, the one statement of the date: \q{Married July 9, 1924.}

The later entries are about the institutions. 5 December 1948, after Lloyd Goodrich has telephoned about the 1950 retrospective: \q{The Whitney has made of itself a clearing house + if one is'nt seen there, no dealer will look at one, it has the effect of an official boycott. \ldots\ She hated women, especially wives.} (p. 56, ref 97076), the \emph{she} being Juliana Force, dead that year; \q{a lone wallflower pinned there to watch everybody else dance.} 19 June 1950, on Goodrich's Penguin monograph: \q{L. Goodrich has seen fit to delete “the artist.” I can never forget that. \ldots\ The Whitney boycott is here complete --- as an artist, historically, I'm not allowed to exist, even in my husband's biography. And that's all right with E.H.!} (p. 61, ref 96058), with a marginal correction that one of two references does say \q{J.N. a painter}. 9 June 1954, a month before the thirtieth anniversary, when he \q{rather generously confessed} that he had not wanted her to paint and now withdrew his opposition, she wrote the anatomy: \q{The 3 points of attack, intent to destroy : 1 - the woman artist. 2. The essential woman + her pretensions to a sex life --- 3. Ability to drive a car, which would include a mind of her own, an intellect.} (p. 73, ref 96115). And a critic's phrase, in 1954, that she found exact: \q{That's why his famous sunlight is without heat. Some critic said that. How extraordinarily perceiving. There is the light on fields + walls of houses, but you'll get no warming benefit --- not from that man} (p. 71, ref 96268).

The book has one sustained happy passage at its end, the group show at the Greenwich Gallery in the spring of 1958, ten canvases: \q{my pictures felt they had waked up in heaven, they having here their prayers, so hopeful to be there where their color sang as never before. I was indeed proud of them.} (p. 77, ref 97417). It breaks off on page 79 with \q{Interrupted} underlined. Forty pages later, with the book turned upside down, two pencil copyings from Milton and Arthur Symons, and nothing else.

Three things about this notebook bear on the ledgers. First, it names the pictures of hers that the ledgers never do: the Obituary with the yellow cat, the Cape Cod bedroom, Washington Square Park from the window at 29 by 36, the studio stove with his crossed feet, Tenements from the End of Williamsburg Bridge finished on 20 February 1928 with the fire escapes left off \q{by both of us}. Second, it explains the paint formula: the woman who wrote \q{Rembrandt colors} under his pictures for six years was a painter who inherited the Rembrandt colours when he changed to Winsor \& Newton in May 1938. Third, and most important for reading Book II, it shows what the naming of figures was. The woman who wrote \q{Shirley} and \q{Olga} and \q{Nora} under his pictures was, in her own book, writing sentences like \q{Dear Miss Chamberlain said last night she loved my houses on tip toe. His houses, she told E.H., she knew had coal in their cellars.} (p. 5). The anecdotes under the paintings are not naive. They are the work of a writer who knew exactly what she was doing and did it in the only book anybody would read.

\section{The Battle of Washington Square}

In August 1946 the trustees of Sailors' Snug Harbor, ground landlords of the north side of Washington Square, gave New York University a twenty-one-year lease on numbers 1, 2 and 3, and the university applied to the Office of Price Administration for certificates of eviction against the tenants. The tenants of number 3 were artists, and the top floor had been Edward Hopper's since December 1913. Two of the four notebooks are about what followed.

\emph{Re: 3 Wash Sq} is a stenographer's spiral pad with \q{Re. 3 Wash. Sq} pencilled on the cover. It is not a diary. The transcription calls it a working draft, and its editorial rule follows from that: \q{The same sentences therefore recur three and four times, and never twice the same: names are added, spelled differently, or dropped. Nothing here is merged.} From the front she wrote drafts: a statement of the case written three times over, the third a fair copy to Whitelaw Reid of the \emph{Herald Tribune} dated 10 February 1947; letters to Mayor LaGuardia in three drafts, to Paul Sachs of the Fogg, to Robert Moses, to Mayor O'Dwyer, to Dean Vanderbilt of the law school, to the rent attorney at the O.P.A.; press lists, guest lists, tax assessments. From the back, and running backwards through the photographs, she wrote a dated chronicle of the fight from January to 25 March, and the two met in the middle of the pad.

The argument of the drafts is about the house. \q{a terrible old house with no heat, high ceilings, long flights to climb, bath rooms - remote but even the bare walls beautiful A house frequently cussed but forever adored.} (\nb{3 Wash Sq}{7}{97317}). \q{Painters require space to paint a picture that will carry across a room or a gallery. \ldots\ He needs north light. Space + light. These 2 things are rare finds - they are too expensive, no new studio buildings are being constructed} (p. 17, ref 97398), marked in red as the point to lead with. \q{The house is specially built for painters --- there are 10 good north light studios + 5 large sky lights. The place reeks with tradition. It is as well adapted to painters as the aquarium to fishes.} (to O'Dwyer, p. 38, ref 96068). The roll of former tenants, Glackens and Edith Dimock, Guy Pène du Bois, Lawson, Rockwell Kent, Dos Passos, Cummings and the \emph{Dial}, Edmund Wilson, Elinor Wylie, Eakins in Stokes's studio, the Gishes, is written out five times and never twice with the same names. \q{The house has produced 9 members of the Institute of Arts + Letters. Can't find one from that N.Y.U.} (p. 14).

And the argument is about him. \q{E. Hopper has just turned the 1/3 mark of a century, having spent the greater part of his life here + painted pictures that hang in museums clean across a continent in this old house} (p. 3, ref 96255). \q{I, as wife, know what it means to have my husband + his 33 years collection sent adrift in an overcrowded city but this issue is really larger than that.} (p. 15). To the O.P.A.: \q{One of the trustees of the Mus. of Mod. Art expressed the fear that if E.H. were so uprooted from \#3 where he has spent the greater part of his life, his work would die. \ldots\ E.H. too is a landmark in Amer. painting of his day.} and, lower down, \q{My husband frowns on my writing you this --- you may or may not find it irrelevant} (pp. 33--34). To Moses: \q{My husband was discouraged over your evasive letter. But my husband is known for understatement. (He leaves out everything that matters, one wonders he speaks up at all.)} (p. 36, ref 97184). The chronicle has him at the first hearing on 18 February, where the university did not appear: \q{E.H. Comment --- since no N.Y.U. there “We tore at our own heartstrings”} (p. 47, ref 97377); at the press party of the 25th, sixty people at the Morses', \q{E.H. said he was for everything that has been said}; and on 18 March, \q{I get letter off to Ira A. Schiller by 11 P.M. E. dissenting plentifully but insisting he should have a woman's point of view.} (p. 44, ref 97538). He paid his share: \q{E.H. received bill from C. Hauptman for 50. + sends check same day. \$50 from each of 6 1/2 tenants} (p. 47).

Her own painting is in the pad at its edges. \q{Took cat water Color drawing to \unc{Teco Gallon} to Juried-in Show at} Macbeth, on 14 February (p. 47); \q{(This 2'' week waiting for visit to see my work by Macbeth Gal.) much depressed.} on 13 March; \q{Louise Morse invited me to hold show in their studio} / \q{E. H. thinks better not.} on the 20th or 21st (p. 44). The fight and the neglect are on the same page, and she did not separate them.

The outcome is the last letter: \q{Beckley phoned 1/2 hr. ago to announce the O.P.A. decision to deny N.Y.U.'s request for certificate of eviction}, 16 May 1947, and \q{The chief value of our victory in this major skirmish is that it reaffirms our profound belief that Wash Sq. belongs to the artists.} (p. 43, ref 97371). The other notebook, \emph{Battle of Wash Sq.}, is a pocket memorandum book of eight sheets that sets the same dates on its first page, \q{Aug. 1946 Sailors Snug gave / lease to N.Y.U.} and \q{May 15th 1947. Decision / from O.P.A.}, and lists the cast (\nb{Battle}{1}{89299}). Its one dated entry is a board meeting of the Washington Square Neighbors on 30 November 1948, at which Owen Grundy of the \emph{Villager} proposed that \q{“residential” be redefined as A - not incl. Educational, hospitals etc. so as to exclude N.Y.U.}, and \q{Jane + I sat back of his shoulder blades.} (p. 3, ref 89300). The rest is a chronicle of demolitions: \q{Rhinelander - Torn down - Dec. 1950 / White Marble House Jan 2'', 51}; \q{S. side of Sq. - Torn down - / Palisade of studio windows / House of Genius}; \q{Summer 1953 - LaFayette demolished}; and \q{Frederick Stokes died at 97 - / Mar. or Feb. 1955}, the dean of the house, whom the drafts had called \q{our polar explorer aged 88} and who \q{climbed up his 74 steps like a mountain goat} (\nb{Battle}{5}{89301}; p. 6). A leaf between pages 8 and 9 was cut out; its stub reads \q{Dec \ldots\ Dec \ldots\ Feb \ldots\ Jan}, a dated run somebody removed.

The victory held for thirteen years and then took a different form. Book III's Whereabouts leaf for 1960 has the university turning five studios into offices and putting in radiators, and the Book I Gifts leaf records the etchings given in May 1947 to the lawyer and the reporter \q{celebrating The Battle of Wash. Sq.} Her Who's Who entry, drafted in the black notebook, put it among her credentials: \q{Active in defense of The Battle of Washington Sq. 1947.} (\nb{Black notebook}{57}{97267}). He died in the house in 1967; the university has it still.

\section{The black notebook}

The black two-ring notebook has nothing written on its cover and no object number. It is, in the transcription's words, \q{a household and studio memorandum book}, and its subject changes every page or two. Its dated content runs from a driver's licence of August 1937 and a paper order of November 1937 to \q{Refound Oct. 14'', 1964.}; the pencil sheets at the back are the oldest stratum and the ink at the front the 1950s. The first page boxes the address, \q{Mrs. Edward Hopper / 3 Washington Square North / New York 3 N.Y. / Phone Spring 7 - 0949. / also Truro, Cape Cod, Mass- / \unc{Her book}} (\nb{Black notebook}{1}{96076}). Her book: the ledgers were his work in her hand; this was hers.

It is the one place in the archive where the costs of the studio are written. Frames, first: the makers, Usui, Heydenryk (spelt three ways), Robinson, Barbizon, Matt of East 54th Street, Burges-Revel of La Porte, Indiana; their prices, \q{Usui frame 25 x 30 - gold, \unc{pinkish} - \$18 / Heydenryk - 25 x 30 - goldish + white. 11.50}; and the one verdict on a frame in the whole corpus, \q{Matt - 25 x 30 was aluminum + gold + no good on anything. Too raw sienna - awful in north light. Gorgeous glittery in fluorescent now worse - more yellow than before.} (p. 18, ref 97168). One bill she ruled as a table, Burges-Revel's of 28 February 1958: model 200 with insert, 29 by 36, finish 2R, 30.75; the same in 25 by 30, 26.75; two of model 214 at 22.75 and 21.00; and a key to what went in them, \q{29 x 36 = Wash. Sq. from Studio window / 25 x 30 = Fire Escape + Convent opposite / 22 x 32 = Buick in Cal. Canyon / 20 x 28 = Cape studio from attic stairs - rug, lamp} (p. 54, ref 97360). These are her pictures. The frames in the black notebook are almost all for her own work, and the one frames page that carries his titles, Chop Suey in \q{old metal} at 32 by 38, Shakespeare at Dusk, is a list of old frames in the rack (p. 122, ref 98173). A recipe for finishing plain wood, taken down from the painter Bertram Hartman, \q{1. Shellac. / 2. White shoe polish liquid / 3. Rub in dry color - pastel - / 4. Wax.} (p. 6, ref 96080), is the one instruction for making anything in a book of lists.

Paper is priced by the pound of the ream and the cent of the sheet: \q{Coper Union - 72 lbs. - 20 ¢ / 140 '' 39 ¢ / See Robinson - 72 '' 23 ¢ / 90 '' 37 ¢ / 140 '' 50 ¢} (p. 9, ref 97564), and the sheet sizes of Whatman and Lindenmeyr are set down for cutting mats. Insurance on the Truro house, a family theft policy of a thousand dollars at \$22.50 and windstorm on \q{Cottage + Garage} at five, taken out in December 1947, stands on a page afterwards cancelled with a pencil cross, the two expiry lines ringed (p. 11, ref 97318). \q{Studio + Bedroom painted 2 Coats - Casein / May 19 1953. - 4 day job. \$250 or 350?} (p. 25). A Macy wrist watch bought in May 1945 for ten dollars, cleaned in 1950 for 3.98, the total \q{7.47} and then \q{7.67 --- Why?} (p. 123). A phonograph, \q{Even exchange for Colombia bgt. Mar 28, 60 / \$153.99}, undelivered and reported to Mrs. Payne at Schirmer's (p. 51). The burglary: \q{1958 - Feb. 11'' Burglar visitation - Money ± 115. + also / Amer. Ex. Checks - ± 900.} (p. 47, ref 97055). A pencil table of his measurements, \q{E. Suit Size}, waist 38, chest 42, \q{Neck - 15 1/2. but work shirts 16.} (p. 16).

Two chronicles of the years are in it, one kept backwards from 1961 to 1906 and one forwards from 1941 to 1958, and they do not meet. The backwards one gives a line to a year and the lines are the skeleton the timeline uses: \q{1934 - Cape - Built house.}; \q{1933 Cape - last yr. at Bird's Cage. / Bgt. land.}; \q{1927 - Two Lights - Coast Guard Cottage \ldots\ got car - Dodge}; \q{1924 Gloucester. Honeymoon. / Watercolors. Chicago August.}; and before 1924 the years are hers alone: \q{1915 - Monhegan ? Wash Sq. Players.}, \q{1914 - Provincetown + Ogunquit}, \q{7 Haarlem, Italy}. The forwards one is the hospitals and the honours: \q{1948. Jan. Group of 6 or 7 paintings at Rehn. / very successful.} then \q{Feb. N.Y. Hospital - Prostate / Grosvenor Hotel when left Hospital / Stayed into May / 3 trips to N.Y. Hos. hemmorhages.}; \q{1950 N.Y. Hospital - Diverticulosis \ldots\ Feb. - Retrospective Ex. Hopper, Man of the Year / at Whitney Mus. - W. 8'' St. E. in hospital / at opening of Exhib.} (p. 45, ref 97471); \q{1957 - Dec. 17 - 1958 Jan. 23'' E.H. to N.Y. Hos. again / Prostate operations - 2 of them.} A page turned sideways puts the Mexican journeys straight, \q{To get this straight / Mexico --- / 1943 Mex. City, Saltillo, Monterrey / 1946 - Sattillo - Grand Tetons / 1951 Sattillo - Santa Fé / 1952 - 53 - El Paso \ldots\ / 1955 \ldots\ / 1956, 57 - Calif.} (p. 48).

And it has the only sentence of Edward Hopper's about his purpose that any of the ten books records, boxed in red crayon:

\begin{leafquote}
Jan 1946. Quoting E. H. / Cruelly pinned down to definition / of his purpose generally, I extract: / “the arrestation of a moment / in time - with the most acute / realization. \hfill (\nb{Black notebook}{64}{98119})
\end{leafquote}

\noindent A second boxed version below it, \q{to arrest / a moment in time / that moment / most acutely realized.}, the transcription is careful to say, is \q{her own recasting of the first, not a second saying of his}. Cruelly pinned down: the phrase is hers, and the word for what she did to get the sentence is the word she used for what he did to her.

The notebook also holds the tenants of number 3 from Glackens on, with \q{E.H. came in 1913} written diagonally beside the list (p. 117, ref 98176); the lions of number 6, \q{S. Kitson facit / John Morron got them 1883 / from Vanderbilt house upper 5'' Ave.} (p. 121); a list of subjects still to paint, \q{Lions \# 6 \ldots\ Eddie at work. / Playground - Wash. Sq. N. / In rear studio - self portrait in “dangerous” hat.} (p. 128, ref 98181); the drawings he gave her, \q{Owned by J.H. gift of E.H. / Helen Hayes House - “Pretty Penny” / preliminary to water color. / \ldots\ Girlie Show - preliminary, crayon.} (p. 59); the two watercolours given to the Museum of Modern Art's soldiers' sale of May 1942 and \q{Not exhibited so took back} (p. 63); a picture \q{stolen at Martha Jackson Gallery 1954} (p. 31); the mailing lists of \emph{Reality}, the realist painters' broadsheet of 1953--55, ninety names with roman numerals for the numbers sent; her Who's Who entry twice; four days of his cold in March 1938 with the temperature taken hour by hour; a reading list that runs from Havelock Ellis to \emph{Ethan Frome}; recipes, \q{Keep rare in middle - so fibre not toughened + unfit to eat.}; and, in 1962, under a heading \q{Little Frying Pans - for girls}, the word \q{Career} written above and struck below (p. 61, ref 97387).

Her marks are a small grammar of their own, and the transcription's glossary records them: a raised double stroke after a numeral is an ordinal, not an inch mark, so \q{Dec. 6'' 1956} is the sixth and \q{46 x 35 1/2''} is inches; \q{c̄} is \emph{with}; \q{±} is \emph{about}; a tick on a mailing list means \q{came or wrote} on one page and \q{wrote or rec'd. by hand} on the next, and both keys are hers. Names are spelt as she spelt them each time, Sattillo and Sortillo and Saltillo, Athenuem and Athenium, and the transcription lets each stand because the leaf did.

\begin{figure}[htbp]\centering
\begin{tikzpicture}
\draw[black!20] (0,0) rectangle (13.00,9.36);
\draw[black!80,line width=0.6pt,-{Stealth[length=3pt]}] (11.21,7.02) to[bend left=8] (11.95,7.55);
\fill[black] (11.21,7.02) circle (1.1pt);
\node[font=\tiny,anchor=east,inner sep=1.5pt] at (11.21,7.02) {Wash. Sq.};
\fill[black] (11.95,7.55) circle (1.1pt);
\node[font=\tiny,anchor=south west,inner sep=1.5pt] at (11.95,7.55) {Gloucester};
\draw[black!80,line width=0.6pt,-{Stealth[length=3pt]}] (11.95,7.55) to[bend left=8] (8.15,7.34);
\fill[black] (8.15,7.34) circle (1.1pt);
\node[font=\tiny,anchor=south,inner sep=1.5pt] at (8.15,7.34) {Chicago};
\draw[black!80,line width=0.6pt,-{Stealth[length=3pt]}] (11.21,7.02) to[bend left=8] (12.05,7.83);
\fill[black] (12.05,7.83) circle (1.1pt);
\node[font=\tiny,anchor=south west,inner sep=1.5pt] at (12.05,7.83) {Two Lights};
\draw[black!80,line width=0.6pt,-{Stealth[length=3pt]}] (12.05,7.83) to[bend left=8] (11.97,7.73);
\fill[black] (11.97,7.73) circle (1.1pt);
\draw[black!80,line width=0.6pt,-{Stealth[length=3pt]}] (11.97,7.73) to[bend left=8] (9.88,4.76);
\fill[black] (9.88,4.76) circle (1.1pt);
\node[font=\tiny,anchor=west,inner sep=1.5pt] at (9.88,4.76) {Charleston};
\draw[black!80,line width=0.6pt,-{Stealth[length=3pt]}] (11.21,7.02) to[bend left=8] (12.09,7.37);
\fill[black] (12.09,7.37) circle (1.1pt);
\node[font=\tiny,anchor=west,inner sep=1.5pt] at (12.09,7.37) {S. Truro};
\draw[black!80,line width=0.6pt,-{Stealth[length=3pt]}] (11.21,7.02) to[bend left=8] (1.98,5.72);
\fill[black] (1.98,5.72) circle (1.1pt);
\node[font=\tiny,anchor=north,inner sep=1.5pt] at (1.98,5.72) {Las Vegas};
\draw[black!80,line width=0.6pt,-{Stealth[length=3pt]}] (1.98,5.72) to[bend left=8] (1.29,5.12);
\fill[black] (1.29,5.12) circle (1.1pt);
\node[font=\tiny,anchor=west,inner sep=1.5pt] at (1.29,5.12) {Los Angeles};
\draw[black!80,line width=0.6pt,-{Stealth[length=3pt]}] (1.29,5.12) to[bend left=8] (0.36,6.18);
\fill[black] (0.36,6.18) circle (1.1pt);
\node[font=\tiny,anchor=west,inner sep=1.5pt] at (0.36,6.18) {San Francisco};
\draw[black!80,line width=0.6pt,-{Stealth[length=3pt]}] (0.36,6.18) to[bend left=8] (0.30,8.38);
\fill[black] (0.30,8.38) circle (1.1pt);
\node[font=\tiny,anchor=west,inner sep=1.5pt] at (0.30,8.38) {Portland};
\draw[black!80,line width=0.6pt,-{Stealth[length=3pt]}] (0.30,8.38) to[bend left=8] (12.09,7.37);
\draw[black!80,line width=0.6pt,-{Stealth[length=3pt]}] (11.21,7.02) to[bend left=8] (5.57,0.94);
\fill[black] (5.57,0.94) circle (1.1pt);
\node[font=\tiny,anchor=east,inner sep=1.5pt] at (5.57,0.94) {Mexico City};
\draw[black!80,line width=0.6pt,-{Stealth[length=3pt]}] (5.57,0.94) to[bend left=8] (5.16,2.67);
\fill[black] (5.16,2.67) circle (1.1pt);
\node[font=\tiny,anchor=east,inner sep=1.5pt] at (5.16,2.67) {Saltillo};
\draw[black!80,line width=0.6pt,-{Stealth[length=3pt]}] (5.16,2.67) to[bend left=8] (5.31,2.75);
\fill[black] (5.31,2.75) circle (1.1pt);
\node[font=\tiny,anchor=east,inner sep=1.5pt] at (5.31,2.75) {Monterrey};
\draw[black!80,line width=0.6pt,-{Stealth[length=3pt]}] (11.21,7.02) to[bend left=8] (5.16,2.67);
\draw[black!80,line width=0.6pt,-{Stealth[length=3pt]}] (5.16,2.67) to[bend left=8] (2.99,7.89);
\fill[black] (2.99,7.89) circle (1.1pt);
\node[font=\tiny,anchor=north,inner sep=1.5pt] at (2.99,7.89) {Grand Teton};
\draw[black!80,line width=0.6pt,-{Stealth[length=3pt]}] (2.99,7.89) to[bend left=8] (12.09,7.37);
\draw[black!80,line width=0.6pt,-{Stealth[length=3pt]}] (11.21,7.02) to[bend left=8] (5.16,2.67);
\draw[black!80,line width=0.6pt,-{Stealth[length=3pt]}] (5.16,2.67) to[bend left=8] (4.05,5.59);
\fill[black] (4.05,5.59) circle (1.1pt);
\node[font=\tiny,anchor=north,inner sep=1.5pt] at (4.05,5.59) {Santa Fe};
\draw[black!80,line width=0.6pt,-{Stealth[length=3pt]}] (4.05,5.59) to[bend left=8] (12.09,7.37);
\draw[black!80,line width=0.6pt,-{Stealth[length=3pt]}] (11.21,7.02) to[bend left=8] (3.93,4.47);
\fill[black] (3.93,4.47) circle (1.1pt);
\node[font=\tiny,anchor=north east,inner sep=1.5pt] at (3.93,4.47) {El Paso};
\draw[black!80,line width=0.6pt,-{Stealth[length=3pt]}] (3.93,4.47) to[bend left=8] (5.16,1.42);
\fill[black] (5.16,1.42) circle (1.1pt);
\node[font=\tiny,anchor=east,inner sep=1.5pt] at (5.16,1.42) {Guanajuato};
\draw[black!80,line width=0.6pt,-{Stealth[length=3pt]}] (5.16,1.42) to[bend left=8] (6.11,0.31);
\fill[black] (6.11,0.31) circle (1.1pt);
\node[font=\tiny,anchor=east,inner sep=1.5pt] at (6.11,0.31) {Oaxaca};
\draw[black!80,line width=0.6pt,-{Stealth[length=3pt]}] (6.11,0.31) to[bend left=8] (5.83,0.80);
\fill[black] (5.83,0.80) circle (1.1pt);
\draw[black!80,line width=0.6pt,-{Stealth[length=3pt]}] (5.83,0.80) to[bend left=8] (5.31,2.75);
\draw[black!80,line width=0.6pt,-{Stealth[length=3pt]}] (5.31,2.75) to[bend left=8] (11.21,7.02);
\draw[black!80,line width=0.6pt,-{Stealth[length=3pt]}] (11.21,7.02) to[bend left=8] (10.44,6.11);
\fill[black] (10.44,6.11) circle (1.1pt);
\node[font=\tiny,anchor=west,inner sep=1.5pt] at (10.44,6.11) {Richmond};
\draw[black!80,line width=0.6pt,-{Stealth[length=3pt]}] (12.09,7.37) to[bend left=8] (11.95,7.55);
\draw[black!80,line width=0.6pt,-{Stealth[length=3pt]}] (11.21,7.02) to[bend left=8] (1.23,5.12);
\fill[black] (1.23,5.12) circle (1.1pt);
\node[font=\tiny,anchor=south west,inner sep=1.5pt] at (1.23,5.12) {Pacific Palisades};
\draw[black!80,line width=0.6pt,-{Stealth[length=3pt]}] (11.21,7.02) to[bend left=8] (11.87,7.48);
\fill[black] (11.87,7.48) circle (1.1pt);
\node[font=\tiny,anchor=east,inner sep=1.5pt] at (11.87,7.48) {Boston};
\draw[black!40,dashed,line width=0.5pt,-{Stealth[length=3pt]}] (11.21,7.02) to[bend left=8] (7.58,6.42);
\fill[black] (11.21,7.02) circle (1.1pt);
\node[font=\tiny,anchor=east,inner sep=1.5pt] at (11.21,7.02) {Wash. Sq.};
\fill[black] (7.58,6.42) circle (1.1pt);
\node[font=\tiny,anchor=south,inner sep=1.5pt] at (7.58,6.42) {St. Louis};
\draw[black!40,dashed,line width=0.5pt,-{Stealth[length=3pt]}] (7.58,6.42) to[bend left=8] (6.59,6.55);
\fill[black] (6.59,6.55) circle (1.1pt);
\draw[black!40,dashed,line width=0.5pt,-{Stealth[length=3pt]}] (6.59,6.55) to[bend left=8] (5.98,6.15);
\fill[black] (5.98,6.15) circle (1.1pt);
\draw[black!40,dashed,line width=0.5pt,-{Stealth[length=3pt]}] (5.98,6.15) to[bend left=8] (4.30,6.48);
\fill[black] (4.30,6.48) circle (1.1pt);
\draw[black!40,dashed,line width=0.5pt,-{Stealth[length=3pt]}] (4.30,6.48) to[bend left=8] (2.72,7.02);
\fill[black] (2.72,7.02) circle (1.1pt);
\draw[black!40,dashed,line width=0.5pt,-{Stealth[length=3pt]}] (2.72,7.02) to[bend left=8] (2.65,6.12);
\fill[black] (2.65,6.12) circle (1.1pt);
\draw[black!40,dashed,line width=0.5pt,-{Stealth[length=3pt]}] (2.65,6.12) to[bend left=8] (2.66,5.69);
\fill[black] (2.66,5.69) circle (1.1pt);
\draw[black!40,dashed,line width=0.5pt,-{Stealth[length=3pt]}] (2.66,5.69) to[bend left=8] (2.46,6.05);
\fill[black] (2.46,6.05) circle (1.1pt);
\draw[black!40,dashed,line width=0.5pt,-{Stealth[length=3pt]}] (2.46,6.05) to[bend left=8] (1.98,5.72);
\fill[black] (1.98,5.72) circle (1.1pt);
\node[font=\tiny,anchor=north,inner sep=1.5pt] at (1.98,5.72) {Las Vegas};
\draw[black!40,dashed,line width=0.5pt,-{Stealth[length=3pt]}] (1.98,5.72) to[bend left=8] (2.10,4.74);
\fill[black] (2.10,4.74) circle (1.1pt);
\draw[black!40,dashed,line width=0.5pt,-{Stealth[length=3pt]}] (2.10,4.74) to[bend left=8] (1.29,5.12);
\fill[black] (1.29,5.12) circle (1.1pt);
\node[font=\tiny,anchor=west,inner sep=1.5pt] at (1.29,5.12) {Los Angeles};
\draw[black!40,dashed,line width=0.5pt,-{Stealth[length=3pt]}] (1.29,5.12) to[bend left=8] (0.96,5.23);
\fill[black] (0.96,5.23) circle (1.1pt);
\draw[black!40,dashed,line width=0.5pt,-{Stealth[length=3pt]}] (0.96,5.23) to[bend left=8] (0.47,5.83);
\fill[black] (0.47,5.83) circle (1.1pt);
\draw[black!40,dashed,line width=0.5pt,-{Stealth[length=3pt]}] (0.47,5.83) to[bend left=8] (0.98,6.17);
\fill[black] (0.98,6.17) circle (1.1pt);
\draw[black!40,dashed,line width=0.5pt,-{Stealth[length=3pt]}] (0.98,6.17) to[bend left=8] (0.37,6.10);
\fill[black] (0.37,6.10) circle (1.1pt);
\draw[black!40,dashed,line width=0.5pt,-{Stealth[length=3pt]}] (0.37,6.10) to[bend left=8] (0.30,8.38);
\fill[black] (0.30,8.38) circle (1.1pt);
\node[font=\tiny,anchor=west,inner sep=1.5pt] at (0.30,8.38) {Portland};
\draw[black!40,dashed,line width=0.5pt,-{Stealth[length=3pt]}] (0.30,8.38) to[bend left=8] (3.01,8.12);
\fill[black] (3.01,8.12) circle (1.1pt);
\node[font=\tiny,anchor=north east,inner sep=1.5pt] at (3.01,8.12) {Yellowstone};
\draw[black!40,dashed,line width=0.5pt,-{Stealth[length=3pt]}] (3.01,8.12) to[bend left=8] (3.26,8.08);
\fill[black] (3.26,8.08) circle (1.1pt);
\draw[black!40,dashed,line width=0.5pt,-{Stealth[length=3pt]}] (3.26,8.08) to[bend left=8] (9.48,7.23);
\fill[black] (9.48,7.23) circle (1.1pt);
\node[font=\tiny,anchor=south,inner sep=1.5pt] at (9.48,7.23) {Cleveland};
\draw[black!40,dashed,line width=0.5pt,-{Stealth[length=3pt]}] (9.48,7.23) to[bend left=8] (11.22,7.12);
\fill[black] (11.22,7.12) circle (1.1pt);
\draw[black!40,dashed,line width=0.5pt,-{Stealth[length=3pt]}] (11.21,7.02) to[bend left=8] (5.72,3.81);
\fill[black] (5.72,3.81) circle (1.1pt);
\draw[black!40,dashed,line width=0.5pt,-{Stealth[length=3pt]}] (5.72,3.81) to[bend left=8] (5.49,3.27);
\fill[black] (5.49,3.27) circle (1.1pt);
\draw[black!40,dashed,line width=0.5pt,-{Stealth[length=3pt]}] (5.49,3.27) to[bend left=8] (5.31,2.75);
\fill[black] (5.31,2.75) circle (1.1pt);
\node[font=\tiny,anchor=east,inner sep=1.5pt] at (5.31,2.75) {Monterrey};
\draw[black!40,dashed,line width=0.5pt,-{Stealth[length=3pt]}] (5.31,2.75) to[bend left=8] (5.16,2.67);
\fill[black] (5.16,2.67) circle (1.1pt);
\node[font=\tiny,anchor=east,inner sep=1.5pt] at (5.16,2.67) {Saltillo};
\draw[black!40,dashed,line width=0.5pt,-{Stealth[length=3pt]}] (5.16,2.67) to[bend left=8] (2.99,7.89);
\fill[black] (2.99,7.89) circle (1.1pt);
\node[font=\tiny,anchor=north,inner sep=1.5pt] at (2.99,7.89) {Grand Teton};
\draw[black!40,dashed,line width=0.5pt,-{Stealth[length=3pt]}] (2.99,7.89) to[bend left=8] (3.01,8.12);
\draw[black!40,dashed,line width=0.5pt,-{Stealth[length=3pt]}] (3.01,8.12) to[bend left=8] (11.21,7.02);
\end{tikzpicture}
\caption{Where the black notebook and the Provincetown typescripts put the Hoppers, drawn at true relative position from coordinates retrieved from OpenStreetMap: solid arrows are Josephine Hopper's year chronicle and her list « To get this straight » of the Mexican journeys, one arrow per leg in the order she wrote them, each year starting from 3 Washington Square North; dashed arrows are the 1941 « Grand Tour of the West » and the 1946 Mexico journey as Madeleine Larson's typescripts give them day by day. It is the order two records fall in, not an itinerary; the ledgers, which record what left the studio and not where anybody was, cannot draw this. After the Timeline tab. Drawn by Claude Fable~5.1, 13 September 2026, from \texttt{travels.json} (2026-09-13); coordinates © OpenStreetMap contributors, ODbL 1.0.}\label{fig:travels}
\end{figure}


\section{What the notebooks do to the ledgers}

Read on their own, the ledgers are the record of a successful painter kept by a devoted wife. Read beside the notebooks, they are the record of two painters kept by the one whose work was not being recorded, in a house they had to fight to keep, with money that arrived through a dealer who died leaving no provision for what he owed, in frames that were \q{awful in north light}. Neither reading is wrong. The ledgers say what a picture cost and where it went; the notebooks say what it cost to be the person who wrote that down. The sentence in the first is the wife's; the sentence in the second is the artist's; and it was the same hand.
